% Options for packages loaded elsewhere
\PassOptionsToPackage{unicode}{hyperref}
\PassOptionsToPackage{hyphens}{url}
\documentclass[
]{article}
\usepackage{xcolor}
\usepackage{amsmath,amssymb}
\setcounter{secnumdepth}{-\maxdimen} % remove section numbering
\usepackage{iftex}
\ifPDFTeX
  \usepackage[T1]{fontenc}
  \usepackage[utf8]{inputenc}
  \usepackage{textcomp} % provide euro and other symbols
\else % if luatex or xetex
  \usepackage{unicode-math} % this also loads fontspec
  \defaultfontfeatures{Scale=MatchLowercase}
  \defaultfontfeatures[\rmfamily]{Ligatures=TeX,Scale=1}
\fi
\usepackage{lmodern}
\ifPDFTeX\else
  % xetex/luatex font selection
\fi
% Use upquote if available, for straight quotes in verbatim environments
\IfFileExists{upquote.sty}{\usepackage{upquote}}{}
\IfFileExists{microtype.sty}{% use microtype if available
  \usepackage[]{microtype}
  \UseMicrotypeSet[protrusion]{basicmath} % disable protrusion for tt fonts
}{}
\makeatletter
\@ifundefined{KOMAClassName}{% if non-KOMA class
  \IfFileExists{parskip.sty}{%
    \usepackage{parskip}
  }{% else
    \setlength{\parindent}{0pt}
    \setlength{\parskip}{6pt plus 2pt minus 1pt}}
}{% if KOMA class
  \KOMAoptions{parskip=half}}
\makeatother
\setlength{\emergencystretch}{3em} % prevent overfull lines
\providecommand{\tightlist}{%
  \setlength{\itemsep}{0pt}\setlength{\parskip}{0pt}}
\usepackage{bookmark}
\IfFileExists{xurl.sty}{\usepackage{xurl}}{} % add URL line breaks if available
\urlstyle{same}
\hypersetup{
  pdftitle={The DIII Winding Number and the Massless Spin(2{,}3) Sector},
  pdfauthor={Aaron Alderman},
  hidelinks,
  pdfcreator={LaTeX via pandoc}}

\title{The DIII Winding Number and the Massless Spin(2,3) Sector}
\author{Aaron Alderman}
\date{18 April 2026}

\begin{document}
\maketitle

\subsection{Abstract}\label{abstract}

We study the topological invariant naturally associated with the
corrected DIII reading of the Spin(2,3) spinor sector. Starting from the
explicit gamma-matrix representation, we construct a chiral
four-generator Clifford frame that anticommutes with the grading
operator \(\Sigma = 2J^{01}\) and defines a flattened chiral Hamiltonian
\(Q(X)\) on \(S^3\). In the chiral basis, \(Q\) has off-diagonal block
\(q(X) \in SU(2)\), and the DIII/AIII winding integral reduces to the
degree of the map \(q : S^3 \to SU(2) \cong S^3\). With the orientation
and sign conventions adopted here, the invariant is \[
W_3 = -1.
\] The overall sign is convention-dependent, but the nontriviality of
the invariant is not. The massless Spin(2,3) sector is therefore best
understood not as a gapped phase carrying the invariant by itself, but
as the equatorial slice inside the same nontrivial Spin(2,3) family.

\subsection{1. Introduction}\label{introduction}

The purpose of this paper is narrow. We do not attempt a complete
topological completion of the Spin(2,3) framework. Instead we isolate
one technical question: once the corrected DIII class assignment is
accepted, what is the corresponding \(d=3\) winding number?

This question matters because the existing topological spine already
fixes most of the algebra. The spinor representation, the grading
operator \(\Sigma\), and the two relevant anti-unitary operators are
already explicit. What remained open was the invariant itself. The aim
here is to compute it in a convention-fixed, line-by-line way and to
state clearly what the result does and does not imply.

The outcome is simple. The natural chiral Spin(2,3) Hamiltonian built
from the Clifford frame defines a map from \(S^3\) into \(SU(2)\), and
that map has degree \(-1\) in the conventions used below. Reversing
orientation or replacing \(q\) by \(q^\dagger\) flips the sign, but not
the fact that the invariant is nonzero. This is the robust content of
the computation.

\subsection{2. Spin(2,3) Conventions and the Corrected Class
Assignment}\label{spin23-conventions-and-the-corrected-class-assignment}

We work with the four-component spinor representation \[
\gamma^0 = i \sigma^2 \otimes \mathbf{1}_2, \qquad
\gamma^1 = i \sigma^1 \otimes \mathbf{1}_2,
\] \[
\gamma^2 = \sigma^3 \otimes \sigma^1, \qquad
\gamma^3 = \sigma^3 \otimes \sigma^2, \qquad
\gamma^4 = \sigma^3 \otimes \sigma^3,
\] with signature \((2,3)\).

The grading operator is taken to be \[
\Sigma = \sigma^3 \otimes \mathbf{1}_2,
\] which agrees with the sector split up to the harmless overall sign
convention relating \(\Sigma\) and \(2J^{01}\).

The relevant anti-unitary symmetries are \[
C = (\gamma^1 \gamma^3)K,
\] and \[
T_0 = -i(\gamma^0 \gamma^3)K,
\] with \[
C^2 = +1, \qquad T_0^2 = -1.
\] The phase of \(T_0\) matters here: the operator is defined only up to
an overall phase, but the above choice is the one for which the chiral
operator is recovered exactly from \(C\) and \(T_0\). Writing the
anti-unitaries as \(C = M_C K\) and \(T_0 = M_{T_0}K\), with \[
M_C = \gamma^1\gamma^3,
\qquad
M_{T_0} = -i\gamma^0\gamma^3,
\] their product is linear and equals \[
CT_0 = M_C M_{T_0}^*.
\] In the present gamma-matrix representation one finds \[
M_C = \sigma^2 \otimes \sigma^2,
\] and \[
M_{T_0}
= -i\gamma^0\gamma^3
=
\begin{pmatrix}
0 & 0 & 0 & 1 \\
0 & 0 & -1 & 0 \\
0 & 1 & 0 & 0 \\
-1 & 0 & 0 & 0
\end{pmatrix},
\] so \[
CT_0
= M_C M_{T_0}^*
=
\begin{pmatrix}
1 & 0 & 0 & 0 \\
0 & 1 & 0 & 0 \\
0 & 0 & -1 & 0 \\
0 & 0 & 0 & -1
\end{pmatrix}
= \Sigma.
\] Thus the chiral operator is not imported from outside the Spin(2,3)
algebra: it is generated internally by the corrected combination of
charge conjugation and partial time reversal. Forgetting the
anti-unitary symmetries yields the chiral unitary class AIII; retaining
them refines the same Hamiltonian family to DIII.

\subsection{3. The Chiral Spin(2,3)
Hamiltonian}\label{the-chiral-spin23-hamiltonian}

Define the following Hermitian matrices: \[
\Gamma_0 = - i\gamma^0 = \sigma^2 \otimes \mathbf{1}_2,
\] \[
\Gamma_1 = \gamma^0\gamma^2 = - \sigma^1 \otimes \sigma^1,
\qquad
\Gamma_2 = \gamma^0\gamma^3 = - \sigma^1 \otimes \sigma^2,
\qquad
\Gamma_3 = \gamma^0\gamma^4 = - \sigma^1 \otimes \sigma^3.
\] These obey \[
\Gamma_a^\dagger = \Gamma_a,
\qquad
\Gamma_a^2 = \mathbf{1}_4,
\qquad
\{\Gamma_a,\Gamma_b\} = 2\delta_{ab}\mathbf{1}_4,
\] for \(a,b = 0,1,2,3\), and each anticommutes with \(\Sigma\): \[
\{\Sigma,\Gamma_a\} = 0.
\]

Hence the Spin(2,3) data contains a natural chiral four-generator
Clifford frame.

Let \[
X = (X_0,X_1,X_2,X_3) \in S^3,
\qquad
X_0^2 + X_1^2 + X_2^2 + X_3^2 = 1,
\] and define \[
Q(X) = X_0 \Gamma_0 + X_1 \Gamma_1 + X_2 \Gamma_2 + X_3 \Gamma_3.
\] Because the \(\Gamma_a\) mutually anticommute and square to the
identity, \[
Q(X)^2 = \mathbf{1}_4,
\] so \(Q\) is already flattened.

This is the correct place where the invariant is defined. The strictly
massless sector is recovered by the equatorial slice \(X_0 = 0\), but
the winding number itself belongs to the full gapped chiral family on
\(S^3\).

\subsection{\texorpdfstring{4. Chiral Basis and the \(SU(2)\)
Block}{4. Chiral Basis and the SU(2) Block}}\label{chiral-basis-and-the-su2-block}

Since \(Q\) anticommutes with \(\Sigma\), it is already off-diagonal in
the chiral basis of \(\Sigma = \sigma^3 \otimes \mathbf{1}_2\). In that
basis one has \[
Q(X)=
\begin{pmatrix}
0 & q_0(X) \\
q_0(X)^\dagger & 0
\end{pmatrix},
\] with \[
q_0(X)= -iX_0\mathbf{1}_2 - (X_1\sigma^1 + X_2\sigma^2 + X_3\sigma^3).
\] This follows directly from the block forms \[
\Gamma_0 =
\begin{pmatrix}
0 & -i\mathbf{1}_2 \\
i\mathbf{1}_2 & 0
\end{pmatrix},
\qquad
\Gamma_j =
\begin{pmatrix}
0 & -\sigma^j \\
-\sigma^j & 0
\end{pmatrix},
\quad j=1,2,3,
\] so the upper-right block of
\(Q(X)=X_0\Gamma_0+X_1\Gamma_1+X_2\Gamma_2+X_3\Gamma_3\) is exactly the
stated \(q_0(X)\).

Now introduce the constant unitary \[
U =
\begin{pmatrix}
e^{-i\pi/4}\mathbf{1}_2 & 0 \\
0 & e^{i\pi/4}\mathbf{1}_2
\end{pmatrix}.
\] Then \[
U^\dagger Q(X) U=
\begin{pmatrix}
0 & e^{i\pi/2}q_0(X) \\
e^{-i\pi/2}q_0(X)^\dagger & 0
\end{pmatrix},
\] and since \(e^{i\pi/2} = i\), the upper-right block becomes \[
q(X)= i q_0(X) = X_0 \mathbf{1}_2 - i(X_1 \sigma^1 + X_2 \sigma^2 + X_3 \sigma^3).
\] Thus after this explicit constant chiral rotation, \(Q\) takes the
standard form \[
Q(X) =
\begin{pmatrix}
0 & q(X) \\
q(X)^\dagger & 0
\end{pmatrix},
\] with \[
q(X) = X_0 \mathbf{1}_2 - i(X_1 \sigma^1 + X_2 \sigma^2 + X_3 \sigma^3).
\] Because \(X \in S^3\), \[
q(X)^\dagger q(X) = \mathbf{1}_2,
\qquad
\det q(X) = 1,
\] so \[
q : S^3 \to SU(2) \cong S^3.
\]

The winding number therefore reduces to the degree of this map. In the
chiral formulation the invariant is \[
W_3
= \frac{1}{24\pi^2}
\int_{S^3}
\operatorname{tr}\!\left[(q^{-1}dq)^3\right].
\]

\subsection{\texorpdfstring{5. Explicit Computation of
\(W_3\)}{5. Explicit Computation of W\_3}}\label{explicit-computation-of-w_3}

Parameterize \(S^3\) by \[
X_0 = \cos\chi,
\qquad
X_1 = \sin\chi \cos\theta,
\] \[
X_2 = \sin\chi \sin\theta \cos\phi,
\qquad
X_3 = \sin\chi \sin\theta \sin\phi,
\] with \[
0 \le \chi \le \pi, \qquad 0 \le \theta \le \pi, \qquad 0 \le \phi < 2\pi.
\]

Write \[
q = X_0 \mathbf{1}_2 - i \vec{X}\cdot\vec{\sigma},
\qquad
q^{-1} = q^\dagger = X_0 \mathbf{1}_2 + i \vec{X}\cdot\vec{\sigma}.
\] Using the Pauli identity \[
\sigma_i \sigma_j = \delta_{ij}\mathbf{1}_2 + i \epsilon_{ijk}\sigma_k,
\] one finds \[
q^{-1}dq
= (X_0\mathbf{1}_2 + iX_i\sigma_i)(dX_0\mathbf{1}_2 - i\,dX_j\sigma_j).
\] Expanding and using \(X_0\,dX_0 + X_i\,dX_i = 0\) on \(S^3\), the
scalar part vanishes and the matrix part reduces to \[
q^{-1}dq
= i\alpha_i \sigma_i,
\] where \[
\alpha_i = X_i\,dX_0 - X_0\,dX_i + \epsilon_{ijk}X_j\,dX_k.
\] Equivalently, defining \(\omega_i = -\alpha_i\), one may write \[
q^{-1}dq = -i \,\omega_i \sigma_i,
\] where the one-forms \(\omega_i\) are the standard left-invariant
coframe on \(SU(2) \cong S^3\).

The cubic trace is then obtained directly from \[
\operatorname{tr}(\sigma_i\sigma_j\sigma_k)=2i\epsilon_{ijk},
\] which gives \[
\operatorname{tr}\!\left[(q^{-1}dq)^3\right]
= 2\,\epsilon_{ijk}\,\alpha_i\wedge\alpha_j\wedge\alpha_k.
\] Substituting the coordinate parameterization of \(S^3\), one obtains
the explicit one-forms \[
\alpha_1 = -\cos\theta\, d\chi + \cos\chi\sin\chi\sin\theta\, d\theta + \sin^2\chi\sin^2\theta\, d\phi,
\] \[
\alpha_2 = -\sin\theta\cos\phi\, d\chi + (-\cos\chi\sin\chi\cos\theta\cos\phi - \sin^2\chi\sin\phi)\, d\theta
+ (\cos\chi\sin\chi\sin\theta\sin\phi - \sin^2\chi\sin\theta\cos\theta\cos\phi)\, d\phi,
\] \[
\alpha_3 = -\sin\theta\sin\phi\, d\chi + (-\cos\chi\sin\chi\cos\theta\sin\phi + \sin^2\chi\cos\phi)\, d\theta
+ (-\cos\chi\sin\chi\sin\theta\cos\phi - \sin^2\chi\sin\theta\cos\theta\sin\phi)\, d\phi.
\] Their coefficient matrix in the basis \((d\chi,d\theta,d\phi)\) has
determinant \[
\det(\alpha_i{}^\mu) = -\sin^2\chi \sin\theta,
\] and hence \[
\epsilon_{ijk}\,\alpha_i\wedge\alpha_j\wedge\alpha_k
= -6 \sin^2\chi \sin\theta \, d\chi \wedge d\theta \wedge d\phi.
\] Therefore \[
\operatorname{tr}\!\left[(q^{-1}dq)^3\right]
= -12 \sin^2\chi \sin\theta \, d\chi \wedge d\theta \wedge d\phi.
\] With the domain orientation taken to be
\(d\chi \wedge d\theta \wedge d\phi\), this is the sign entering the
winding integral in the present convention. This coordinate orientation
is the standard outward orientation on \(S^3\): indeed \[
\det\!\bigl(X,\partial_\chi X,\partial_\theta X,\partial_\phi X\bigr)
= \sin^2\chi \sin\theta > 0
\] away from the coordinate singular sets. The minus sign in \(W_3\)
therefore does not come from the parameterization of \(S^3\); it comes
from the winding-number convention used here, namely \[
W_3 = \frac{1}{24\pi^2}\int_{S^3}\operatorname{tr}\!\left[(q^{-1}dq)^3\right].
\] Many references instead insert the opposite overall sign, or
equivalently use the opposite Maurer-Cartan convention, in which case
the same map \(q(X)=X_0\mathbf{1}_2-i\vec X\cdot\vec\sigma\) would be
assigned \(+1\).

Therefore \[
W_3
= \frac{1}{24\pi^2}
\int_0^\pi \int_0^\pi \int_0^{2\pi}
(-12)\sin^2\chi \sin\theta \, d\phi\, d\theta\, d\chi.
\] Evaluating the integrals, \[
\int_0^\pi \sin^2\chi\, d\chi = \frac{\pi}{2},
\qquad
\int_0^\pi \sin\theta\, d\theta = 2,
\qquad
\int_0^{2\pi} d\phi = 2\pi,
\] so \[
W_3
= \frac{-12}{24\pi^2}
\left(\frac{\pi}{2}\right)(2)(2\pi)
= -1.
\]

Thus the topological invariant is nontrivial. Its overall sign depends
on orientation and on the sign convention in the definition of \(q\);
replacing \(q\) by \(q^\dagger\), or reversing orientation on \(S^3\),
yields \(+1\) instead. The convention-independent statement is
therefore:

\begin{quote}
the Spin(2,3) chiral Hamiltonian carries \(|W_3| = 1\).
\end{quote}

\subsection{6. Interpretation of the Massless
Sector}\label{interpretation-of-the-massless-sector}

The computation above applies to the gapped chiral family \(Q(X)\) on
\(S^3\). The strictly massless sector corresponds to the equatorial
slice \[
X_0 = 0,
\] for which \[
q = - i(X_1 \sigma^1 + X_2 \sigma^2 + X_3 \sigma^3).
\] This equator is not itself a gapped three-dimensional DIII phase
carrying \(W_3\). The computation above therefore does not by itself
prove a full interface theorem. What it does show is that the massless
slice sits naturally as the equator of a gapped chiral family with
nontrivial unit winding. It is therefore reasonable to regard it as a
candidate transition locus within that family, but that last
interpretive step goes beyond the bare invariant calculation.

The relation to the earlier AIII language is then straightforward. If
one suppresses the anti-unitary symmetries and keeps only the chiral
grading, the same flattened Hamiltonian lies in class AIII. Including
\(T_0\) and \(C\) refines that same chiral Hamiltonian to class DIII.
The integer computed here is therefore the same winding number under
both descriptions, even though AIII and DIII classify different
symmetry-constrained Hamiltonian families. It is most precise to present
the result in DIII language because the corrected Spin(2,3) symmetry
data contains those anti-unitaries explicitly.

\subsection{7. What the Result Does and Does Not
Establish}\label{what-the-result-does-and-does-not-establish}

The result established here is exact and narrow. The corrected Spin(2,3)
symmetry data yields a natural chiral Clifford frame, that frame defines
a flattened Hamiltonian on \(S^3\), and the associated winding number is
nontrivial, equaling \(-1\) in the present orientation convention. This
is the firm mathematical content of the paper.

What is not established is equally important. The calculation does not
by itself provide a material realization of the corresponding DIII
Hamiltonian, an anomaly-inflow statement, a generation-counting theorem
from topology, or a full extension of the analysis to the broader
two-time-reversal picture. Those questions remain open and should be
treated separately from the invariant computed here.

\subsection{8. Numerical Cross-Check}\label{numerical-cross-check}

A numerical integration of the winding formula using a direct
discretization of the \(S^3\) coordinates reproduces the analytic
result. With the current script and a uniform \(28 \times 28 \times 28\)
grid in \((\chi,\theta,\phi)\), one finds \[
W_3 \approx -0.999.
\] The residual error is discretization error and decreases under grid
refinement. A script implementing this check is included with the source
files for this paper.

\subsection{9. Conclusion}\label{conclusion}

The corrected topological Spin(2,3) story now has a concrete invariant.
The natural chiral Hamiltonian extracted from the explicit gamma-matrix
representation defines a map \(S^3 \to SU(2)\) of degree one in
magnitude, giving \(W_3 = -1\) in the conventions adopted here. The
strictly massless sector is therefore best viewed as the equatorial
slice of a gapped chiral family with nontrivial unit winding, rather
than as an isolated gapped phase carrying the invariant on its own.
Interpreting that slice as a genuine transition locus is plausible, but
it remains one step beyond what the winding-number calculation alone
proves.

\end{document}
